% GENERATED FILE. Edit canonical lessons, then run npm run generate:books.
% canonical-source-hash: fnv1a64:5923c6b8b7570d16
% canonical-lessons: RU-C327-arkheolog, RU-C327-raskopki, RU-C327-nakhodka, RU-C327-sokrovishche, RU-C327-klad

\chapter{Treasure and finds}
\label{ch:ru-treasure-and-finds}
\hlchaptermodality{327}

\begin{chapteropening}
\textbf{By the end of this chapter:} \emph{I can talk about an archaeologist, a dig, a find, treasure and a buried hoard.}

Complete the last lesson of chapter 327: I can talk about an archaeologist, a dig, a find, treasure and a buried hoard.
\end{chapteropening}

\section[\texorpdfstring{\ru{археолог} (arkheológ)}{археолог (arkheológ)}]{\ru{археолог} (arkheológ) — an archaeologist}

\label{lesson:RU-C327-arkheolog}



\begin{quote}
Before the new one: say the Russian for a prince of old Rus, then the Russian for to found.
\end{quote}

\subsection*{You'll want to know: \ru{археолог}}
\textbf{\ru{археолог}} (\emph{arkheológ}) — \textquotedblleft{}an archaeologist\textquotedblright{}.

\begin{quote}
\ru{Археолог} \ru{работал} \ru{в} \ru{пустыне} \ru{два} \ru{года}.
\end{quote}

\begin{quote}
\emph{Arkheológ rabótal v pustýne dva góda.}
\end{quote}

\begin{quote}
The archaeologist worked in the desert for two years.
\end{quote}

\emph{\ru{Археология}} — archaeology.

\begin{cousinweb}
\emph{\ru{археолог}} is built on Greek \emph{arkhaîos} \textquotedblleft{}ancient\textquotedblright{}, from \emph{arkh\'{\={e}}} \textquotedblleft{}beginning\textquotedblright{}, and \emph{lógos} \textquotedblleft{}word, account\textquotedblright{} (EtymDB): one who studies ancient things.
\end{cousinweb}

\begin{grammarlens}[title={What you've built}]
One more word for telling, explaining and describing at B1.
\end{grammarlens}

\subsection*{Your turn}
\begin{itemize}
\raggedright
  \item \emph{Say it:} \emph{arkheológ}
  \item \emph{Say it:} \emph{arkheológ}, once more
  \item \emph{Say it:} a sentence of your own with \emph{arkheológ}
  \item \emph{Recall:} say the Russian for a prince of old Rus, then the Russian for to found, then say \emph{arkheológ} again
\end{itemize}

\begin{tcolorbox}[breakable,colback=teal!4,colframe=teal!35!black,title={Before you move on}]
What does \ru{археолог} mean? (\textquotedblleft{}An archaeologist\textquotedblright{}.) Say it once more.
\end{tcolorbox}
\section[\texorpdfstring{\ru{раскопки} (raskópki)}{раскопки (raskópki)}]{\ru{раскопки} (raskópki) — a dig, an excavation}

\label{lesson:RU-C327-raskopki}



\begin{quote}
Before the new one: say the Russian for to found, then the Russian for an archaeologist.
\end{quote}

\subsection*{You'll want to know: \ru{раскопки}}
\textbf{\ru{раскопки}} (\emph{raskópki}) — \textquotedblleft{}a dig, an excavation\textquotedblright{}.

\begin{quote}
\ru{Летом} \ru{студенты} \ru{ездили} \ru{на} \ru{раскопки}.
\end{quote}

\begin{quote}
\emph{Létom studénty yézdili na raskópki.}
\end{quote}

\begin{quote}
In the summer the students went on a dig.
\end{quote}

Plural; genitive \emph{\ru{раскопок}}.

\begin{cousinweb}
\emph{\ru{раскопки}} is built on \emph{\ru{раскопать}} \textquotedblleft{}to dig up\textquotedblright{}, from \emph{\ru{рас}-} and \emph{\ru{копать}} \textquotedblleft{}to dig\textquotedblright{}, inherited from Proto-Slavic \emph{kopati} (EtymDB).
\end{cousinweb}

\begin{grammarlens}[title={What you've built}]
One more word for telling, explaining and describing at B1.
\end{grammarlens}

\subsection*{Your turn}
\begin{itemize}
\raggedright
  \item \emph{Say it:} \emph{raskópki}
  \item \emph{Say it:} \emph{raskópki}, once more
  \item \emph{Say it:} a sentence of your own with \emph{raskópki}
  \item \emph{Recall:} say the Russian for to found, then the Russian for an archaeologist, then say \emph{raskópki} again
\end{itemize}

\begin{tcolorbox}[breakable,colback=teal!4,colframe=teal!35!black,title={Before you move on}]
What does \ru{раскопки} mean? (\textquotedblleft{}A dig, an excavation\textquotedblright{}.) Say it once more.
\end{tcolorbox}
\section[\texorpdfstring{\ru{находка} (nakhódka)}{находка (nakhódka)}]{\ru{находка} (nakhódka) — a find, a discovery}

\label{lesson:RU-C327-nakhodka}



\begin{quote}
Before the new one: say the Russian for an archaeologist, then the Russian for a dig.
\end{quote}

\subsection*{You'll want to know: \ru{находка}}
\textbf{\ru{находка}} (\emph{nakhódka}) — \textquotedblleft{}a find, a discovery\textquotedblright{}.

\begin{quote}
\ru{Эта} \ru{находка} \ru{удивила} \ru{учёных}.
\end{quote}

\begin{quote}
\emph{Éta nakhódka udivíla uchyónykh.}
\end{quote}

\begin{quote}
This find amazed the scientists.
\end{quote}

Feminine; genitive plural \emph{\ru{находок}}.

\begin{cousinweb}
\emph{\ru{находка}} (EtymDB: \textquotedblleft{}a find, a discovery\textquotedblright{}) is built on \emph{\ru{находить}} \textquotedblleft{}to find\textquotedblright{}, from \emph{\ru{на}-} and \emph{\ru{ходить}} \textquotedblleft{}to go\textquotedblright{}, Proto-Slavic \emph{xoditi} (EtymDB): what you come upon.
\end{cousinweb}

\begin{grammarlens}[title={What you've built}]
One more word for telling, explaining and describing at B1.
\end{grammarlens}

\subsection*{Your turn}
\begin{itemize}
\raggedright
  \item \emph{Say it:} \emph{nakhódka}
  \item \emph{Say it:} \emph{nakhódka}, once more
  \item \emph{Say it:} a sentence of your own with \emph{nakhódka}
  \item \emph{Recall:} say the Russian for an archaeologist, then the Russian for a dig, then say \emph{nakhódka} again
\end{itemize}

\begin{tcolorbox}[breakable,colback=teal!4,colframe=teal!35!black,title={Before you move on}]
What does \ru{находка} mean? (\textquotedblleft{}A find, a discovery\textquotedblright{}.) Say it once more.
\end{tcolorbox}
\section[\texorpdfstring{\ru{сокровище} (sokróvishche)}{сокровище (sokróvishche)}]{\ru{сокровище} (sokróvishche) — a treasure}

\label{lesson:RU-C327-sokrovishche}



\begin{quote}
Before the new one: say the Russian for a dig, then the Russian for a find.
\end{quote}

\subsection*{You'll want to know: \ru{сокровище}}
\textbf{\ru{сокровище}} (\emph{sokróvishche}) — \textquotedblleft{}a treasure\textquotedblright{}.

\begin{quote}
\ru{Пираты} \ru{спрятали} \ru{сокровище} \ru{на} \ru{острове}.
\end{quote}

\begin{quote}
\emph{Piráty spryátali sokróvishche na óstrove.}
\end{quote}

\begin{quote}
The pirates hid the treasure on an island.
\end{quote}

Neuter; often plural \emph{\ru{сокровища}}.

\begin{cousinweb}
\emph{\ru{сокровище}} is usually traced to \emph{\ru{сокровенный}} \textquotedblleft{}hidden\textquotedblright{}, built on \emph{\ru{со}-} and the root of \emph{\ru{крыть}} \textquotedblleft{}to cover\textquotedblright{}: what is kept hidden. EtymDB lists no parent.
\end{cousinweb}

\begin{grammarlens}[title={What you've built}]
One more word for telling, explaining and describing at B1.
\end{grammarlens}

\subsection*{Your turn}
\begin{itemize}
\raggedright
  \item \emph{Say it:} \emph{sokróvishche}
  \item \emph{Say it:} \emph{sokróvishche}, once more
  \item \emph{Say it:} a sentence of your own with \emph{sokróvishche}
  \item \emph{Recall:} say the Russian for a dig, then the Russian for a find, then say \emph{sokróvishche} again
\end{itemize}

\begin{tcolorbox}[breakable,colback=teal!4,colframe=teal!35!black,title={Before you move on}]
What does \ru{сокровище} mean? (\textquotedblleft{}A treasure\textquotedblright{}.) Say it once more.
\end{tcolorbox}
\section[\texorpdfstring{\ru{клад} (klad)}{клад (klad)}]{\ru{клад} (klad) — a buried treasure, a hoard}

\label{lesson:RU-C327-klad}



\begin{quote}
Before the new one: say the Russian for a find, then the Russian for a treasure.
\end{quote}

\subsection*{You'll want to know: \ru{клад}}
\textbf{\ru{клад}} (\emph{klad}) — \textquotedblleft{}a buried treasure, a hoard\textquotedblright{}.

\begin{quote}
\ru{В} \ru{саду} \ru{старого} \ru{дома} \ru{нашли} \ru{клад}.
\end{quote}

\begin{quote}
\emph{V sadú stárovo dóma nashlí klad.}
\end{quote}

\begin{quote}
A hoard was found in the garden of an old house.
\end{quote}

Masculine.

\begin{cousinweb}
\emph{\ru{клад}} is built on the root of \emph{\ru{класть}} \textquotedblleft{}to put, to lay\textquotedblright{}, which EtymDB traces to PIE \emph{kleh\textsubscript{2}-} \textquotedblleft{}to put, to lay out\textquotedblright{}: what was laid away.
\end{cousinweb}

\begin{grammarlens}[title={What you've built}]
One more word for telling, explaining and describing at B1.
\end{grammarlens}

\subsection*{Your turn}
\begin{itemize}
\raggedright
  \item \emph{Say it:} \emph{klad}
  \item \emph{Say it:} \emph{klad}, once more
  \item \emph{Say it:} a sentence of your own with \emph{klad}
  \item \emph{Recall:} say the Russian for a find, then the Russian for a treasure, then say \emph{klad} again
\end{itemize}

\begin{tcolorbox}[breakable,colback=teal!4,colframe=teal!35!black,title={Before you move on}]
What does \ru{клад} mean? (\textquotedblleft{}A buried treasure, a hoard\textquotedblright{}.) Say it once more.
\end{tcolorbox}
